\chapter{Gene regulatory networks}

\section{Regulation of gene expression}
\subsection{Overview: coordinate the gene interactions}

\subsection{Regulation of the chromatin structure}

\subsection{Regulation of transcription initiation}

\subsection{Post-transcriptional regulation}


\section{Tools to model biological networks}
Instead of considering each gene as independent and static, the analyses reported in this chapter aim at recovering the holistic and dynamic structure of the transcriptome.

\section{Overview: network design and paradigms}

\paragraph{Top-down vs. Bottom-up}
\paragraph{Knowledge vs. Pure Learning approaches}
\todo[fancyline]{refer to the Boolean network internship}
\paragraph{Model graphs}

\begin{enumerate}[label=\alph*)]
\item[Boolean networks:]
\item[Random networks:]
\item[Mechanistic networks:]

\end{enumerate}

\section{Gaussian graphical models}
\subsection{Conditional independence in multivariate Gaussian distributions}


\subsection{Learn \glsentryshort {ggm} in high-dimensional setting}


\subsection{From \glsentryshort {ggm} to \glsentryshort{gbn}}

\paragraph{Overview of \glsentryshort {gbn}}
\paragraph{Learn a \gls{gbn} under topological constraints}

\todo[fancyline]{\cite{dahl_etal05}}
\subsection{Differential network analysis}

\paragraph{Overview: dynamic identification of key biological drivers}
Predict dynamically the behaviour of a new drug and prioritise biomarkers integrating the graph structure
\subsubsection{A local-score approach for the identification of key biological drivers}




